\chapter{Conclusion and Future Work}
\label{ch:conclusion}

\section{Conclusion}
RetailPulse set out to turn the footage a store already records into service and staffing information on an ordinary CPU, without identifying anyone. At the end of FYP-1 the core loop works end to end on recorded and synthetic data: the vision pipeline finds people, labels customers and employees, and confirms interactions with a pose check; each run becomes validated events in PostgreSQL; two business agents answer manager questions with figures that are checked against their source; and an unattended customer becomes a task for an eligible employee, without a manager approval step and never twice for the same event.

Measured against the objectives of Chapter~\ref{ch:intro}:
\begin{itemize}
\item \textbf{O1, roles:} all 56 matched held-out tracks classified correctly with the fitted weights.
\item \textbf{O2, interactions:} F1 0.450 on held-out clips (interval 0.26--0.63); the pose check raises precision from 0.346 to 0.391 at equal recall.
\item \textbf{O3, event store:} every run's rows pass the contract check; a synthetic day covers all event tables.
\item \textbf{O4, questions:} 51 golden questions are routed, answered and refused correctly, with every answer fully grounded (after fixes made on the same set).
\item \textbf{O5, dispatch:} dispatch, no eligible employee, duplicates and blind cameras behave correctly in tests and in the synthetic rehearsal.
\item \textbf{O6, CPU only:} 0.60 times real time on an idle CPU.
\item \textbf{O7, access:} the security review's findings are fixed and tested.
\end{itemize}

The mid evaluation found the report's figures and requirements incorrect. This report rebuilds both: the requirements specification is now numbered, atomic and testable with measurable non-functional targets, and every design figure has been redrawn from the current code.

The main limitation is evidence, not function: no acceptance scenario has yet passed on real store footage, the backend has not been deployed, and the labelled set is small enough that the accuracy interval is wide.

\section{FYP-2 Plan}
\begin{longtable}{P{0.27\textwidth}P{0.45\textwidth}P{0.2\textwidth}}
\caption{FYP-2 work plan}\label{tab:fyp2}\\
\toprule
Work & Done when & Owner\\
\midrule
\endfirsthead
\toprule
Work & Done when & Owner\\
\midrule
\endhead
\bottomrule
\endfoot
Real-footage acceptance & A1--A10 pass on recorded store footage through the integrated system & Ruhab (lead), all\\
Store map and calibration & Floor plan and camera points supplied and registered; A2 runs; map view shows real cells & Saqib, Moiz\\
Backend deployment & Backend reachable from the deployed web tier with production secrets; S1 checked in the deployment & Saqib\\
Automatic ingestion & A run started from the web page is stored without a manual command & Saqib\\
Second annotator and blind re-check & Tuning labels second-reviewed; IMG\_8838 part~2 and IMG\_8830 re-checked & Labeller, Ruhab\\
Dispatch agreement & Human rater's blind choices compared with the system on all 20 scenarios & Ruhab\\
Clip review on real events & Two reviewers; precision and Cohen's kappa reported (A10) & Ruhab, Saqib\\
Task transitions & Backend rejects out-of-order task states & Saqib\\
Camera handoff & Labelled crossings collected and the handoff evaluator scored & Saqib\\
Branch protection & CI required before merging to the main branch & Moiz\\
Model and agent freeze & Weights and configuration frozen and recorded before the final demonstration & All\\
\end{longtable}

\section{Future Work Beyond FYP-2}
Cross-camera re-identification without biometric identity, metric floor distances with calibrated static cameras, GPU deployment for real-time processing, shelf-level product interaction, and operation across several stores remain outside the project scope. They are recorded so a later project can start from the same event store and agent contracts.
